\chapter{The equation as a beginning, not a promise}\label{ch:beginning}
\question{What can we now say precisely, and what remains open?}

Return to the clinic before its first patient arrives. The pump, the payment record and the need for safe care are still there. We have not replaced them with a formula. We have learned to ask what part of the event a formula can describe and what work is required before its number deserves an interpretation.

A physical action has a declared boundary and measurable quantities. A feasible action respects the constraints of the represented world. A reference and scale supply an explicit ruler. A potential measures the declared condition. A marginal describes a local sensitivity. The finite account compares complete endpoints and includes separately justified residual process burden once. A policy then uses or ignores that information under its own rules. Actors choose; execution and audit establish what occurred.

\section{The clinic revisited}
In the teaching world, delivering three litres from $(14,6)$ to $(11,9)$ reduced the declared potential by $3.75$. We can explain that number, reproduce it and show why a larger delivery can have a different sign. We can also explain why the unchanged surroundings cancel within an additive model.

We cannot infer from it that the clinic is safe, every patient is served or a real pump used no energy. Those would require more state variables, measurements and evidence. Nor does the positive number identify which institution should own a receipt, who pays the electricity bill or what a technician should earn.

Suppose the pump needs repair. The repair's immediate represented physical effects belong in its current account. If subsequent measurement establishes a changed capacity or efficiency, certification can support a new model version for future actions. Old records retain the version under which they were made. The problem of identifying and certifying persistent improvement remains a domain task; the equation does not pay immediately for an unsupported forecast.

An institution must separately authorise the work, carry any assigned settlement responsibility and compensate the worker. A negative current physical account is not personal blame. A beneficial future field change does not prove that the same worker will capture all its benefits. Public provision and access guarantees therefore remain substantive institutional questions.

Distance deserves the same discipline. A nearby service may require less transport, but proximity has no automatic EBU coefficient. A distant route may be more efficient under a particular physical account. The comparison must follow represented quantities, losses, time and capacity rather than an unexamined preference for either local or distant provision.

\section{The proposed experiment, briefly}
\textit{Prospective model proposal, not an executed result.} The Local Gaussian programme would specify external physical forcing; a frozen state; EBU-blind candidate actions or groups; joint physical feasibility; exact valuation and declared attributions; owner-specific capacity affordability; unweighted random actor choice; exact selected execution; and settlement and audit. No hidden value optimiser, service-first fallback or restoring dynamics supplies success in advance.\cite{gaussian}

A physical-random comparison would use declared physical rules without the capacity-affordability restriction. The future protocol must define action menus, owners, empty choices, numerical rules, disturbances and outcome measures before inspecting results. Sharing random draws would not automatically ensure equal applied disturbances after states diverge and feasibility differs.

The account might help, make little difference or perform poorly under that study. An introductory book cannot choose among those outcomes. It can make the question clear enough that the study is capable of answering it. No parameters, seeds, launch plan or experimental implementation are supplied by this edition.

\section{Where the other books belong}
The preserved Part II remains the home of rigorous mathematics under its named models. Its threshold potentials and controller proofs keep their original scope. A Gaussian specialisation is an additional declared construction, not a retroactive change to those proofs.

Part III preserves implementation and evidence history. D0, P1C and other historical findings must be read with their own designs and manifests. They are not Gaussian experimental evidence. This new volume changes the route into the subject without rewriting that record.

The current series architecture gives later volumes distinct responsibilities. Part IV concerns measurement, evidence and falsification. Part V concerns recovery and stochastic dynamics. Part VI develops multiple actions, interaction and local structure. Part VII follows routes, transformations and boundary effects across distance. Part VIII studies optional coordination, memory and feedback mechanisms. Part IX addresses institutions, rights, access, responsibility and transition.\cite{series}

Those responsibilities are research tasks, not guaranteed chapters of favourable results. In particular, a proof of accounting closure does not remove the need for causal identification, an acceptable ownership rule or a defensible treatment of people who cannot earn experimental capacity.

\section{What the reader can now challenge}
A careful reader can ask whether an action is physically defined, whether quantity and rate have been confused, whether the reference is compatible and whether it is reachable. They can check the sign convention, demand a finite endpoint, identify a missing affected factor and refuse to count the same burden twice.

They can also ask whether a plotted curve is an equation, an observation or a simulation. They can distinguish an unchanged total from a recovered function and a path allocation from a right to payment. These habits make the proposal more testable. They do not require faith in its eventual success.

The central question remains whether this explicit account is useful in worlds that can be measured and institutions that people can justify. The evidence of ecological pressures and unequal access makes that question worth investigating. The evidence needed to answer it is still separate.

We finish with the account whose meaning we can now state:
\begin{equation*}
 \boxed{\Delta E_a=V_A(z)-V_A(T_a(z))-C_a.}
\end{equation*}
One declared baseline. One coherent physical transformation. Complete affected support. A fixed, justified ruler. Compatible units and no double counting. The calculation is a beginning: EBU calculates, actors choose, and evidence must establish what follows.
